%%%%%%%%%%%%%%%%%%%%%%%%%%%%%%%%%%%%%%%%%%%%%%%
% include additional packages you need to use
%%%%%%%%%%%%%%%%%%%%%%%%%%%%%%%%%%%%%%%%%%%%%%%
\usepackage{appendix}
\usepackage{url}			% for \url in bibliography
\usepackage{adjustbox}		% 표 폭이 컬럼을 넘으면 자동 축소 (max width)

% graphic, float package
\usepackage{graphicx}		% for setting images
\usepackage{float}			% for float objects
\usepackage{subfloat}		
\usepackage{subfigure}		% for adding several figures in a figure environment
\usepackage{lscape}			% for landscape type images or tables

\usepackage{enumitem}

% for compact section title spacing
% \usepackage[compact]{titlesec}


% mathmetical presentation
\usepackage{gensymb}
\usepackage{amsmath}
\usepackage{amssymb}
\usepackage{amsthm}
\usepackage{exscale}
\usepackage{textcomp}		% extra symbols


% for circled number
\newcommand{\cl}[1]{\textcircled{\scriptsize #1}}


% package for using algorithmic presentation (설치된 경우에만 로드)
\IfFileExists{algorithm.sty}{%
  \usepackage{algorithmic}%
  \usepackage{algorithm}%
  \renewcommand{\algorithmicrequire}{\makebox[40px]{\hfill\textbf{Input :}}}%
  \renewcommand{\algorithmicensure}{\makebox[40px]{\hfill\textbf{Output :}}}%
}{}

% array and table presentation
\usepackage{array}
\usepackage{tabulary}
\usepackage{multirow}
\usepackage[table]{xcolor}
\usepackage{tikz}
\usetikzlibrary{patterns,positioning,fit}
\usepackage{pgfplots}
\pgfplotsset{compat=1.18}
\usepackage{quantikz}		% 양자회로(ansatz) 다이어그램 (xcolor 뒤에 로드)
\usepackage{ctable}
\usepackage{booktabs}		% for typesetting tables at the level of publication
\usepackage{refcount}		% \getrefnumber (번호 발음 기반 자동 조사)
